\hyphenation{уст-рой-ства}  % помогаем latex понять, как разделять это слова по слогам


\section*{Цель работы}
Изучить принцип работы синхронных двоичных датчиков с использованием среды Verilog, проанализировать их RTL интерпретации и временные диаграммы.

% === ВАЖНО:
% •	Все временные диаграммы должны быть построены ТОЛЬКО С УЧЕТОМ ЗАДЕРЖКИ! Диаграммы без учета задержки в отчете не нужны.
% •	Чтобы на временных диаграммах подписи сигналов и цифры внутренних состояний счетчиков не были слишком мелкими, при формировании скриншотов временных диаграмм не рекомендуется разворачивать окно с построенными диаграммами на весь экран.
% •	Сигналы на временных диаграммах в пп. 1–4 должны быть расположены в следующем порядке. Первым должен быть тактовый сигнал. Предпоследним – выходной сигнал. Последним – сигнал переноса (если он фигурирует в коде. Если нет, последним должен быть выходной сигнал). Остальные сигналы (при наличии) расположить в произвольном порядке.
% •	На временных диаграммах в пп. 1–4 в выходном сигнале должны быть видны ВСЕ РАЗРЯДЫ ВСЕХ ВНУТРЕННИХ СОСТОЯНИЙ! По умолчанию внутренние состояния отображаются в двоичном виде, однако если в столбце слева (где перечислены имена сигналов) выбрать выходной сигнал, нажать правую кнопку мыши, выбрать Radix и затем выбрать Unsigned Decimal, внутренние состояния будут отображаться в десятичном виде. Подобная смена позволит снизить число разрядов числа, отображаемого в шестиугольнике, обозначающем внутреннее состояние счетчика.
% •	На временных диаграммах в пп. 1–2 должны быть показаны все внутренние состояния счетчика, а также сброс в ноль по достижении счетчиком максимального значения. Время моделирования подбирать в соответствии с этим (после сброса в 0 дальше строить не нужно).

% •	На временных диаграммах в пп. 3–4 в виде меандра должен быть задан только тактовый сигнал. Другие входные одноразрядные сигналы задавать в виде меандров не нужно. Кроме того, ВХОДНЫЕ СИНХРОННЫЕ СИГНАЛЫ НЕ ДОЛЖНЫ ИЗМЕНЯТЬСЯ В МОМЕНТЫ РАБОЧИХ ПЕРЕПАДОВ ТАКТОВОГО СИГНАЛА!
% •	На временных диаграммах в п. 3:
% 	Должна быть показана работа счетчика в обоих режимах работы.
% 	Обязательно должны быть показаны моменты срабатывания сигналов переноса и заема.
% 	Не нужно задавать слишком много значений сигналу данных. Двух-трёх вполне будет достаточно.
% 	Не нужно делать больше двух загрузок одного значения сигнала данных.
% 	Сигнал загрузки не должен действовать дольше, чем два такта подряд.

% •	На временных диаграммах в п. 4 необходимо показать следующие режимы работы счетчика:
% 	Счёт есть;
% 	Счёта нет;
% 	Действие сигнала сброса при разрешении счёта;
% 	Действие сигнала сброса при отсутствии счёта.


\section*{ \fbox{ Задание 1. } }
\begin{quote}
    {\itshape
        Собрать схему синхронного двоичного четырехразрядного счетчика (см. \cref{lst:code1.v}).
        
        Изучить схему, реализованную в RTL Viewer.

        Построить временные диаграммы, иллюстрирующие работу устройства. Период тактового сигнала задать 35 нс.

        При каком внутреннем состоянии счетчика, согласно теории, сигнал переноса должен быть равен единице? При каком внутреннем состоянии счетчика сигнал переноса равен единице на временных диаграммах? Если эти два состояния не совпадают, измените код 1 так, чтобы установка сигнала переноса в единицу происходила согласно теории.

        Построить схему в RTL Viewer и временные диаграммы для исправленной версии кода (если исправление проводилось).
    }
\end{quote}

\begin{codemultipage}
    \captionof{listing}{Код 1. Синхронный двоичный четырехразрядный счетчик.\label{lst:code1.v}}
    \nopagebreak
    \inputminted{verilog}{scripts/code1.v}
\end{codemultipage}

RTL схема и временная диаграмма представлены на \cref{fig:code1_RTL.pdf,fig:code1_vwf.png}.

\begin{figure}[H]
    \centering
    \includegraphics[width=1\textwidth]{figs/code1_RTL.pdf}
    \caption{RTL. Код 1.}
    \label{fig:code1_RTL.pdf}
\end{figure}
\begin{figure}[H]
    \centering
    \includegraphics[width=1\textwidth]{figs/code1_vwf.png}
    \caption{Временная диаграмма. Код 1.}
    \label{fig:code1_vwf.png}
\end{figure}

\subsection*{Анализ RTL схемы}

Перед анализом \cref{fig:code1_RTL.pdf} следует ознакомиться с новым элементом \mintinline{text}§Add0§. Он суммирует два входных сигнала. В случае Кода 1 такими являются константа \mintinline{text}§4'h1§ (где \mintinline{text}§h§ означет hex --- шестнадцатиричную систему счисления) и значение \mintinline{text}§out§. С элементами \mintinline{text}§...~reg0§ мы уже ознакомились в предыдущих лабораторных работах --- это D-триггеры.

Итак, к анализу схемы. В момент рабочего такта в триггер \mintinline{text}§out§ записывается значение $x_n = x_{n-1} + 1$ благодаря сумматору. Получившийся результат с выхода \mintinline{text}§Q§ передается в \mintinline{text}§4and§ и на выход \mintinline{text}§out§. Значение после блока И передается на вход триггера \mintinline{text}§p_out§, а затем на выход \mintinline{text}§p_out§.

\subsection*{Проверка теории}

Согласно теории, сигнал переноса должен быть равен единице в момент, когда внутреннее состояние счетчика равно максимальному его значению. В данном случае максимальное значение счетчика равно $1111_2=[2^4-1]_{10}=15_{10}$. Однако на \cref{fig:code1_vwf.png} хорошо видно, что сигнал переноса истиннен, когда значение счетчика равно $0$.

Объяснить это можно как через RTL схему (в рамках данного, не слишком сложного, Кода 1), так и через Verilog код:
\begin{itemize}
    \item RTL. В момент рабочего такта сигнал подается в один и тот же момент времени на оба D-триггера. Оба триггера обновляют свои значения в один момент времени. Однако в \mintinline{text}§p_out§ таким образом записывается необновленное значение \mintinline{text}§out§.
    \item Verilog. И \mintinline{text}§p_out§, и \mintinline{text}§out§ являются переменными, поскольку объявлены через \mintinline{text}§reg§. Обновление их состояний происходит только в рамках блока \mintinline{text}§always§ с соотвествующем условием в списке чувствительности. Имея опыт в языках программирования, изученных в прошлых семестрах (Python, C, C++, MatLab), кажется, что все правильно: сначала обновляется состояние \mintinline{text}§out§ и только затем \mintinline{text}§p_out§. Однако здесь заключается особенность языка Verilog. Он является не столько языком программирования, сколько языком описания схем. Поэтому последовательность строк 12--13 не имеет значения --- состояние обоих переменных обновляется одновременно в момент времени \mintinline{text}§posedge clk§ и поэтому данная реализация счетчика не соответствует теории.
\end{itemize}

\subsection*{Корректная реализация счетчика}

Для того, чтобы значение сигнала переноса корректно изменялось с нуля на единицу в то время, когда счетчик равен $15$, следует использовать оператор непрерывного присваивания \mintinline{text}§assign§. Его особенность заключается в том, что работает он постоянно; в отличие от \mintinline{text}§reg§, который изменяется только в рамках блока \mintinline{text}§always§, срабатывающий в конкретные моменты времени. Исправленный Код 1 представлен на \cref{lst:code1_fixed.v}.

\begin{codemultipage}
    \captionof{listing}{Исправленный Код 1.\label{lst:code1_fixed.v}}
    \nopagebreak
    \inputminted{verilog}{scripts/code1_fixed.v}
\end{codemultipage}

Покажем также временную диаграмму на \cref{fig:code1_fixed_vwf.png}.
\begin{figure}[H]
    \centering
    \includegraphics[width=1\textwidth]{figs/code1_fixed_vwf.png}
    \caption{Временная диаграмма исправленного Кода 1.}
    \label{fig:code1_fixed_vwf.png}
\end{figure}

Примечательно, что сигнал \mintinline{text}§p_out§ не определен на временных диаграммах, а его настоящее значение показывается на сигнале \mintinline{text}§p_out[0]§. Многовероятно, что это связано с определением \mintinline{text}§output [0:0] p_out§, что эквивалентно \mintinline{text}§output p_out§. Он не был убран с диаграмм для сохранения прозрачности работы Quartus II.


\section*{ \fbox{ Задание 2. } }
\begin{quote}
    {\itshape
        Модифицировать схему синхронного двоичного четырехразрядного счетчика так, чтобы получился двоично-десятичный счетчик. Для этого воспользоваться кодом из \cref{lst:code2.v}.

        Изучить схему, реализованную в RTL Viewer.

        Построить временные диаграммы, иллюстрирующие работу устройства. Период тактового сигнала задать 25 нс.

        Что означает выражение \mintinline{verilog}§if (out<4'd9)§? Как изменится работа счетчика, если в данной строке кода число 9 заменить на число 10?
    }
\end{quote}

\begin{codemultipage}
    \captionof{listing}{Код 2. Двоично-десятичный счетчик.\label{lst:code2.v}}
    \nopagebreak
    \inputminted{verilog}{scripts/code2.v}
\end{codemultipage}

RTL схема и временная диаграмма представлены на \cref{fig:code2_RTL.pdf,fig:code2_vwf.png}.

\begin{figure}[H]
    \centering
    \includegraphics[width=1\textwidth]{figs/code2_RTL.pdf}
    \caption{RTL. Код 2.}
    \label{fig:code2_RTL.pdf}
\end{figure}
\begin{figure}[H]
    \centering
    \includegraphics[width=1\textwidth]{figs/code2_vwf.png}
    \caption{Временная диаграмма. Код 2.}
    \label{fig:code2_vwf.png}
\end{figure}

\subsection*{Анализ RTL схемы}
Перед анализом \cref{fig:code2_RTL.pdf} следует ознакомиться с новым элементом \mintinline{text}§LESS_THAN§. Данный элемент получает на вход сигналы A и B и выполняет следующую проверку:
\begin{equation}
    LESS\_THAN = 
    \begin{cases}
        1 & \text{если } A  <  B, \\
        0 & \text{если } A \ge B.
    \end{cases}
    \label{eq:LT}
\end{equation}

С остальными элементами, включая мультиплексор \mintinline{text}§MUX21§ мы уже знакомы.

Итак, к анализу схемы. В момент рабочего такта одновременно обновляются состояния триггеров \mintinline{text}§out§ и \mintinline{text}§p_out§. Для простоты с этого момента первый будем обозначать как $O$, второй --- как $P$. Аналогично Заданию 1, в данной реализации счетчика значение $P$ отстает, что подтверждается временной диаграммой на \cref{fig:code2_vwf.png} --- $P=1$ когда $O=0$.

Выходной сигнал $Q$ триггера $O$ подается на вход сумматору и сравнивателю так, что, если значение $O>9$, срабатывает мультиплексор, обнуляющий сигнал $O$. В ином случае, если $O\le9$, мультиплексор на входе $SEL$ получает ноль, в качестве выходного сигнала передается значение со входа $DATAB$.

Таким образом, данная реализация счетчика аналогична прошлой, за исключением, что в данном случае его обнуление происходит только если $O>9$ (несмотря на то, что он способен хранить значение до $2^4-1=15$), в то время как в счетчике из Задания 1 обнуление происходило простым переполнением.

\subsection*{Анализ выражения}

Выражение \mintinline{verilog}§if (out<4'd9)§ реализует сравниватель \cref{eq:LT}. Согласно проанализировнной RTL схеме в пункте выше при ложности данного неравенства $LESS\_THAN(A=Q_O, B=9)=0$ выполняется обнуление $O=0$, не давая счетчику считать больше установленного значения в 4 бита, представленного десятичной девяткой: $\overline{abcd}_2=9_{10}$.

Если заменить $9$ на $10$, то сравниватель \cref{eq:LT} будет срабатывать при достижении счетчиком $O$ значения $11$ и обнулять его. Счетчик будет считать в диапазоне $0..10$.

\section*{ \fbox{ Задание 3. } }
\begin{quote}
    {\itshape
        Собрать схему синхронного двоичного трёхразрядного реверсивного счетчика с синхронной загрузкой данных. Для этого воспользоваться кодом из \cref{lst:code3.v}.

        Изучить схему, реализованную в RTL Viewer.

        Построить временные диаграммы, иллюстрирующие работу устройства. Период тактового сигнала задать 40 нс.

        В чём принципиальное отличие в формировании сигнала переноса по сравнению со счетчиками из пп. 1-–2? Как это отличие отражается на временных диаграммах?
    }
\end{quote}
% Чтобы при построении временных диаграмм задать какое-либо значение сигналу, представленному в виде шины, в окне Simulation Waveform Editor необходимо выделить произвольный фрагмент сигнала на временной оси (точно так же, как выделяют одноразрядный двоичный сигнал, чтобы установить его в состояние логической единицы или логического нуля) и нажать кнопку Arbitrary Value. В открывшемся окне установить желаемое значение сигнала в строке Numeric or named value с учетом системы счисления, регулируемой с помощью выпадающего списка Radix. Нажать ОК.
\begin{codemultipage}
    \captionof{listing}{Код 3. Синхронный двоичный трёхразрядный реверсивный счетчик с синхронной загрузкой данных.\label{lst:code3.v}}
    \nopagebreak
    \inputminted{verilog}{scripts/code3.v}
\end{codemultipage}

RTL схема и временные диаграммы представлены на \cref{fig:code3_RTL.pdf,fig:code3_vwf_expanded.png,fig:code3_vwf_collapsed.png}.

\begin{figure}[H]
    \centering
    \includegraphics[width=1\textwidth]{figs/code3_RTL.pdf}
    \caption{RTL. Код 2.}
    \label{fig:code3_RTL.pdf}
\end{figure}
\begin{figure}[H]
    \centering
    \includegraphics[width=1\textwidth]{figs/code3_vwf_expanded.png}
    \caption{Развернутая временная диаграмма. Код 2.}
    \label{fig:code3_vwf_expanded.png}
\end{figure}
\begin{figure}[H]
    \centering
    \includegraphics[width=1\textwidth]{figs/code3_vwf_collapsed.png}
    \caption{Сжатая временная диаграмма. Код 2.}
    \label{fig:code3_vwf_collapsed.png}
\end{figure}

\subsection*{Анализ RTL схемы}
Новых элементов на \cref{fig:code3_RTL.pdf} нет.

Рассматрим $P$. Из схемы видно, что для сигнала переноса $P$ больше нет своего триггера --- теперь этот сигнал формируется путем последовательного применения операция И, ИЛИ, НЕ. Из RTL схемы довольно затруднительно корректно считать его логику, поэтому обратимся к Verilog коду: 

\mintinline{verilog}§assign p_out = up & (&out) | ~up & ~(|out)§. 

У данного счетчика есть сигнал $UP$: если $UP=1$, то счетчик инкрементируется, иначе декреметируется. В то же время $P=1$ только если счетчик достиг своей конечной величины. Это объясняет причину такой непростой, если сравнивать с предыдущими счетчиками, логикой. Суть же ее заключается следующем: если счетчик инкрементирует $+1$, то конечная величина $111_2$ должна состоять из всех единиц; если же счетчик декрементирует $-1$, то конечная величина $000_2$ должна быть из всех нулей. Описанная логика для \mintinline{text}§p_out§ реализует данную проверку побитовыми И, ИЛИ, НЕ.

На схеме так же присутствуют два мультиплексора. Один из них (тот, что правее) передает в зависимости от сигнала $load$ на вход триггеру $O$ либо значение $data$, либо значение левого мультиплексора. Логика же левого мультиплексора нам уже известна из прошлых Заданий, но с корректировкой на сигнал $up$. Если $up=1$, то значение $O$ должно увеличиваться с каждым тактом. Тогда левый мультиплексор передает сигнал со входа $DATAB$. В случае $up=0$ сигнал передается со входа $DATAA$. Таким образом становится ясно, что сумматор $Add1$ выполняет $x_n = x_{n-1} - 1$, а $Add0$ --- $x_n = x_{n-1} + 1$. 

\subsection*{Принципиальное отличие}
Принципиальное отличие сигнала переноса в данной реализации от предыдущих заключается в использовании \mintinline{text}§assign§ вместо \mintinline{text}§reg§. Детальная логика работы $P$ описана в подразделе выше, а отличие \mintinline{text}§assign§ от \mintinline{text}§reg§ в Задании 1, разделе <<Корректная реализация счетчика>>.

\section*{ \fbox{ Задание 4. } }
\begin{quote}
    {\itshape
        Собрать схему синхронного двоичного четырехразрядного счетчика с сигналами сброса и разрешения работы. Для этого воспользоваться кодом из \cref{lst:code4.v}.
        
        Синхронными или асинхронными являются сигналы сброса и разрешения работы? Объяснить, почему. Ответ не должен быть основан на временных диаграммах, они являются следствием, а не причиной.

        Изучить схему, реализованную в RTL Viewer.

        Построить временные диаграммы, иллюстрирующие работу устройства. Период тактового сигнала задать 20 нс.
    }
\end{quote}

\begin{codemultipage}
    \captionof{listing}{Код 4. Синхронный двоичный четырехразрядный счетчик с сигналами сброса и разрешения работы.\label{lst:code4.v}}
    \nopagebreak
    \inputminted{verilog}{scripts/code4.v}
\end{codemultipage}

RTL схема и временная диаграмма представлены на \cref{fig:code4_RTL.pdf,fig:code4_vwf.png}.

\begin{figure}[H]
    \centering
    \includegraphics[width=1\textwidth]{figs/code4_RTL.pdf}
    \caption{RTL. Код 2.}
    \label{fig:code4_RTL.pdf}
\end{figure}
\begin{figure}[H]
    \centering
    \includegraphics[width=1\textwidth]{figs/code4_vwf.png}
    \caption{Временная диаграмма. Код 2.}
    \label{fig:code4_vwf.png}
\end{figure}

\subsection*{(А)синхронность}

Сигналы сброса $reset$ и разрешения работы $enable$ являются синхроннымы, поскольку их логика описана внутри блока \mintinline{text}§always§, который срабатывает только во время переднего фронта тактов \mintinline{text}§posedge clk§.

Асинхронными же называются сигналы, значение которые считывается сразу после их изменения. 

\section*{ \fbox{ Задание 5. } }
\begin{quote}
    {\itshape
        Собрать схему синхронного двоичного 32-разрядного счетчика согласно коду из \cref{lst:code5.v}.

        Изучить схему, реализованную в RTL Viewer. Временное моделирование не проводить.

        Если входная частота счётчика составляет 50 МГц, а линии выходной шины \mintinline{verilog}§q[28:25]§ подключены к светодиодам LED5--LED2, то с какой частотой мигает светодиод LED5? Поясните почему.
    }
\end{quote}

\begin{codemultipage}
    \captionof{listing}{Код 5. Синхронный двоичный 32-разрядный счетчик.\label{lst:code5.v}}
    \nopagebreak
    \inputminted{verilog}{scripts/code5.v}
\end{codemultipage}

RTL схема представлена на \cref{fig:code5_RTL.pdf}.

\begin{figure}[H]
    \centering
    \includegraphics[width=1\textwidth]{figs/code5_RTL.pdf}
    \caption{RTL. Код 2.}
    \label{fig:code5_RTL.pdf}
\end{figure}

\subsection*{Анализ RTL схемы}
Новых элементов на \cref{fig:code5_RTL.pdf} нет.

Особенность данного счетчика заключается в его 32-х битность, в остальном нам уже известна логика данной RTL схемы, опираясь на разборы выше.

Триггер счетчика $O$ обновляет свое значение только по переднему фроную тактового сигнала, суммируя предыдущее значение с единицей. Подпись у сумматора подтверждает это: $32'h0000\_0001$. $32$ означает число битов, $h$ указывает систему счисления (в данном случае шестнадцатиричная). У 16-ричной системы счисления одна цифра может принимать значения $0..F$, что соответствует $0..15$ в десятиричной системе счисления. Таким образом одна цифра в 16-ричной системе счисления --- это 4 цифры в двоичной системе счисления ($15_{10}=1111_2$). Всего цифр у входа $B$ сумматора $4\cdot2=8 \implies 8\cdot4=32$. Все сходится, ровно 32 бита.

\subsection*{Светодиоды}
Если частота счетчика $50$ МГц, значит самый младший разряд 32-х битного числа меняется 50 миллионов раз в секунду. Второй справа разряд будет изменяться уже 25 миллионов раз в секунду, поскольку ему нужно <<подождать>>, пока младший разряд дважды изменит свое значние. Причем мерцание --- это $0 \rightarrow 1 \rightarrow 0$, а значит $частота = \frac{кол-во~изменений}{2}$. Таким образом
\begin{equation}
    f_i
    = \frac{
        f_0
    }{
        2^{i + 1}
    }.
    \label{eq:fi}
\end{equation}

$28$ бит подается на $LED5$, $27$ --- на $LED4$, $26$ --- на $LED3$, $25$ --- на $LED2$. По стандартам Quartus II первый бит справа имеет индекс 0 --- формула \cref{eq:fi}  соответствует этому стандарту. Тогда
\begin{equation}
    f_{28}
    = \frac{
        50~МГц   
    }{
        2^{28 + 1}
    }
    = \pv[-3.2]{50e6 / 2**(28 + 1)}~мГц.
\end{equation}


\section*{Выводы}

% === Отчет должен содержать:
% - Тексты программ;
% - Скриншоты схем в RTL Viewer;
% - Временные диаграммы;
% - Ответы на вопросы из каждого пункта
% - Выводы


В ходе выполнения лабораторной работы были проанализированы RTL схемы и временные диаграммы счетчиков.

Было выяснено, что для корректного считывания сигнала переноса $P$ необходимо использовать \mintinline{text}§assign§. Использование \mintinline{text}§reg§ внутри блока \mintinline{text}§always§ приводит к запаздыванию на один такт.

Классический двоичный счетчик обнуляется за счет переполнения. Для реализации счетчиков с произвольным пределом требуется внедрение дополнительного опрератора \mintinline{text}§if-else§, который осуществляет сброс значения $O$ при достижении заданного порогового значения.

Сигналы сброса $reset$ и разрешения работы $enable$ в исследуемых схемах являются синхронными, так как их логика обрабатывается внутри блока \mintinline{text}§always @(posedge clk)§. Схемы реагируют на уровни данных сигналов только в момент переднего фронта тактового сигнала.

Доказано, что в многоразрядном двоичном счетчике легко выполнять деление частоты. Частота изменения каждого разряда подчиняется $f_i = \cfrac{f_{0}}{2^{i+1}}$.